\chapter{Properties of the Fourier Transform and the Convolution Theorem}

In the last few chapters we built the Fourier transform from the ground up: a
signal is a vector, sinusoids are an orthogonal basis, and the transform is
the list of inner products of the signal with each pure
frequency. That geometric picture -- \emph{projection onto rotating phasors} --
forms the foundation for everything in this chapter. Every operational property
we establish is a direct geometric consequence of how rotating vectors behave.

The danger with this material is that it is usually taught as a table to be
memorized: ``time shift becomes a phase ramp, scaling flips the width,
convolution becomes multiplication.'' A memorized table is easy to
forget. Once you understand it, each line becomes a short geometric statement
about what happens to a bundle of spinning vectors when you nudge, stretch, or smear the
signal. We are going to \emph{see} each property, then prove it in two or three
lines, then use it.

One property, the convolution theorem, underlies much of digital signal
processing. It is why filtering can be fast, why the FFT is so useful, and why a
microcontroller on your wrist can do in microseconds what looks like a heavy
integral. We give it the most attention.

\begin{focusbox}
By the end of this chapter you should be able to look at a manipulation of a
signal (delay it, speed it up, multiply it by a cosine, slide one signal across
another) and \emph{predict} what happens to its spectrum without computing a
single integral -- and you should see clearly why convolution in
time becomes multiplication in frequency.
\end{focusbox}

\section{The conventions we will use}
\label{ch11:conventions}

Fourier conventions are a minefield of factors of $2\pi$. We fix one and stick
to it. For a continuous signal $x(t)$ we write the transform in terms of
ordinary frequency $f$ (in hertz):
\begin{align}
X(f) &= \FT\{x\}(f) = \int_{-\infty}^{\infty} x(t)\, e^{-\jj 2\pi f t}\, \dd t,
\label{ch11:ft-fwd}\\
x(t) &= \int_{-\infty}^{\infty} X(f)\, e^{+\jj 2\pi f t}\, \dd f.
\label{ch11:ft-inv}
\end{align}
This is the cleanest convention: no stray $2\pi$ out front, and the forward and
inverse transforms are symmetric. We will occasionally switch to angular
frequency $\omega = 2\pi f$ when it makes a formula prettier, and we will note
when we do. We write the transform-pair relationship as
\[
x(t) \;\longleftrightarrow\; X(f).
\]

For discrete signals $x[n]$ the discrete-time Fourier transform (DTFT) and the
DFT have the same properties with $\int$ replaced by $\sum$ and a periodic
frequency axis; we will state the discrete versions where they differ in a way
that bites you in practice (the DFT's shifts are \emph{circular}, which matters
enormously and we will not let you forget it).

\begin{intuitionbox}
Hold one mental image throughout: $X(f)$ is a \emph{recipe of phasors}. The
number $X(f)$ is a complex amplitude -- a length and a starting angle -- telling
you how much of the unit-speed rotation $e^{\jj 2\pi f t}$ to add in, and at what
phase, to reconstruct $x(t)$. Every property below is the answer to ``what does
this operation do to all those lengths and starting angles at once?''
\end{intuitionbox}

\section{Linearity: the transform is a linear map}
\label{ch11:linearity}

The first property is the one we already used to \emph{define} the transform, so
it almost does not feel like a property. If $a$ and $b$ are constants,
\[
a\,x(t) + b\,y(t) \;\longleftrightarrow\; a\,X(f) + b\,Y(f).
\]
The proof is one line: the integral in \eqref{ch11:ft-fwd} is linear, so it
passes through sums and scalars. Done.

Linearity is simple, but a great deal depends on it. It is what lets us \emph{build
signals out of pieces and reason about each piece separately}. When you see a
recorded forearm tremor as ``a slow drift plus a 10\,Hz physiological tremor
plus broadband sensor noise,'' linearity tells you that the
spectrum of the sum is the sum of the spectra -- you can find and remove the
drift in the frequency domain without disturbing the tremor.

\begin{whybox}
Linearity is the statement that the Fourier transform is a \emph{change
of basis} (Chapter~3). A change of basis is a linear map by
construction: rewriting a vector in new coordinates cannot care whether you
scaled it or added two vectors before or after rewriting. Every linear-algebra
fact about coordinate changes -- and there are many -- is now available to you
for free.
\end{whybox}

\begin{connectionbox}
In the linear-algebra companion, a linear map between finite-dimensional spaces
is a matrix. The DFT is literally a matrix: $X = F x$ where $F_{kn} =
e^{-\jj 2\pi kn/N}$. Linearity of the transform is just the statement
$F(ax+by) = aFx + bFy$, which is true of \emph{every} matrix. The Fourier
matrix happens to be a special one (unitary up to scale, with rows that are
the orthogonal phasor basis), but linearity asks nothing special of it.
\end{connectionbox}

\section{Time shift: delay becomes a phase ramp}
\label{ch11:time-shift}

Suppose we delay the signal by $t_0$
seconds, forming $x(t - t_0)$. What happens to its spectrum?

\begin{align}
\FT\{x(t-t_0)\}(f)
&= \int_{-\infty}^{\infty} x(t-t_0)\, e^{-\jj 2\pi f t}\, \dd t \nonumber\\
&= \int_{-\infty}^{\infty} x(u)\, e^{-\jj 2\pi f (u+t_0)}\, \dd u
\qquad (u = t - t_0) \nonumber\\
&= e^{-\jj 2\pi f t_0} \int_{-\infty}^{\infty} x(u)\, e^{-\jj 2\pi f u}\, \dd u
\nonumber\\
&= e^{-\jj 2\pi f t_0}\, X(f).
\label{ch11:time-shift-result}
\end{align}

So
\[
x(t - t_0) \;\longleftrightarrow\; e^{-\jj 2\pi f t_0}\, X(f).
\]
The \emph{magnitude} of the spectrum is completely unchanged:
$\abs{e^{-\jj 2\pi f t_0}} = 1$. Only the phase changes, and it changes by a
\emph{linear ramp} in frequency: $-2\pi f t_0$.

\begin{happeningbox}
Why a \emph{ramp}, and why proportional to $f$? Each frequency component is a
phasor spinning at $f$ revolutions per second. To delay the whole signal by
$t_0$, you must rewind every phasor to where it was $t_0$ seconds ago. A fast
phasor (large $f$) covered more angle in $t_0$ seconds than a slow one, so it has
to be rewound by more angle. The angle to rewind is exactly
$(\text{rev/s}) \times (\text{s}) \times 2\pi = 2\pi f t_0$. Faster phasors get
more phase shift -- and ``proportional to $f$'' is the definition of a straight
line through the origin. A pure delay is a \emph{linear phase} shift. The idea
returns when we design filters that must not smear the timing of
a gesture.
\end{happeningbox}

\begin{trapbox}
A constant delay gives \emph{linear} phase. The reverse is the lesson designers
care about (Figure~\ref{ch11:fig-linear-phase}): if a system's phase response is
\emph{not} a straight line versus frequency, then different frequencies are delayed by different amounts, and a
sharp event (a tap, a finger flick) gets smeared in time. Linear phase $=$ pure
delay $=$ no waveform distortion. Nonlinear phase $=$ frequency-dependent delay
$=$ smearing. This is why ``linear-phase FIR filters'' are prized in
gesture and audio work.
\end{trapbox}

\begin{figure}[H]
\centering
\begin{tikzpicture}
\begin{axis}[
  width=12cm, height=4.2cm,
  xlabel={frequency $f$}, ylabel={phase $\angle$ (rad)},
  xmin=0, xmax=6, ymin=-9.5, ymax=0.5,
  xtick={0,1,2,3,4,5,6}, ytick={0,-2,-4,-6,-8},
  axis lines=left, samples=2, domain=0:6,
  every axis plot/.append style={thick}, clip=true,
]
\addplot[sigblue] {-1.4*x};
\node[sigblue] at (axis cs:4.2,-3.2) {\small slope $=-2\pi t_0$};
\end{axis}
\end{tikzpicture}
\caption{Time shift adds a phase ramp. A pure delay $t_0$ multiplies the spectrum
by $e^{-\jj 2\pi f t_0}$, leaving magnitude untouched and tilting the phase along
a straight line of slope $-2\pi t_0$. A larger delay tilts the line more
steeply.}
\label{ch11:fig-linear-phase}
\end{figure}

\section{Frequency shift (modulation): multiply by a phasor, slide the spectrum}
\label{ch11:freq-shift}

By Fourier symmetry, time shift has a mirror image. Multiply the
signal by a complex phasor $e^{\jj 2\pi f_0 t}$:
\begin{align}
\FT\{e^{\jj 2\pi f_0 t}\, x(t)\}(f)
&= \int_{-\infty}^{\infty} x(t)\, e^{\jj 2\pi f_0 t}\, e^{-\jj 2\pi f t}\, \dd t
\nonumber\\
&= \int_{-\infty}^{\infty} x(t)\, e^{-\jj 2\pi (f - f_0) t}\, \dd t
= X(f - f_0).
\label{ch11:freq-shift-result}
\end{align}
Multiplying by a phasor in time \emph{slides the entire spectrum} up by $f_0$:
\[
e^{\jj 2\pi f_0 t}\, x(t) \;\longleftrightarrow\; X(f - f_0).
\]

This is the exact dual of the time-shift property -- delay in time is a phase
ramp in frequency; a phase ramp in time (which is what $e^{\jj 2\pi f_0 t}$ is)
is a shift in frequency. The two formulas are the same statement viewed from the
two sides of the transform.

\subsection{Real modulation and the cosine}

Multiplying by a real cosine is the practical case. Using
$\cos(2\pi f_0 t) = \tfrac12(e^{\jj 2\pi f_0 t} + e^{-\jj 2\pi f_0 t})$ and
linearity:
\[
\cos(2\pi f_0 t)\, x(t) \;\longleftrightarrow\;
\tfrac12 X(f - f_0) + \tfrac12 X(f + f_0).
\]
The spectrum is \emph{copied} -- one copy slid up to $+f_0$, one slid down to
$-f_0$, each at half height. This is amplitude modulation, and the picture
(Figure~\ref{ch11:fig-modulation}) is worth memorizing.

\begin{figure}[H]
\centering
\begin{tikzpicture}
\begin{axis}[
  width=12cm, height=4.2cm,
  xlabel={frequency $f$}, ylabel={$\abs{\cdot}$}, ylabel style={rotate=-90},
  xmin=-9, xmax=9, ymin=0, ymax=1.35,
  xtick={-6,-3,0,3,6}, xticklabels={$-f_0$,,$0$,,$f_0$},
  ytick=\empty, axis lines=left, samples=120, domain=-9:9,
  every axis plot/.append style={thick}, clip=true,
]
\addplot[sigblue] {max(0, 1 - abs(x)/3)};
\node[sigblue] at (axis cs:0,1.15) {\small $X(f)$};
\end{axis}
\end{tikzpicture}

\begin{tikzpicture}
\begin{axis}[
  width=12cm, height=4.2cm,
  xlabel={frequency $f$}, ylabel={$\abs{\cdot}$}, ylabel style={rotate=-90},
  xmin=-9, xmax=9, ymin=0, ymax=1.35,
  xtick={-6,-3,0,3,6}, xticklabels={$-f_0$,,$0$,,$f_0$},
  ytick=\empty, axis lines=left, samples=120, domain=-9:9,
  every axis plot/.append style={thick}, clip=true,
]
\addplot[specred] {0.5*max(0, 1 - abs(x-6)/3)};
\addplot[specred] {0.5*max(0, 1 - abs(x+6)/3)};
\node[specred] at (axis cs:6,0.72) {\small $\tfrac12 X(f-f_0)$};
\node[specred] at (axis cs:-6,0.72) {\small $\tfrac12 X(f+f_0)$};
\end{axis}
\end{tikzpicture}
\caption{Modulation by a cosine. Top: the baseband spectrum $X(f)$ of a slowly
varying signal. Bottom: after multiplying by $\cos(2\pi f_0 t)$, the spectrum is
copied to $\pm f_0$ at half height. The shape is preserved; it is only
\emph{relocated}.}
\label{ch11:fig-modulation}
\end{figure}

\begin{applicationbox}
This is how AM radio works, how an SDR tunes a station, and how I/Q
demodulation pulls a signal down to baseband (Chapter~23). It
is also how a sensor's spectrum gets corrupted: if a motor or a switching power
supply injects a tone at $f_0$ into your wrist IMU's supply rail, any real
fluctuation $x(t)$ riding on that rail appears in your data as a copy of $X(f)$
\emph{moved up to} $f_0$. Recognizing modulation in a spectrum -- ``two
symmetric copies of the same shape'' -- is a useful diagnostic clue when debugging hardware.
\end{applicationbox}

\section{Scaling: stretch in time, squeeze in frequency}
\label{ch11:scaling}

Speed the signal up or slow it down by a factor $a > 0$, forming $x(at)$.
\begin{align}
\FT\{x(at)\}(f)
&= \int_{-\infty}^{\infty} x(at)\, e^{-\jj 2\pi f t}\, \dd t \nonumber\\
&= \frac{1}{a}\int_{-\infty}^{\infty} x(u)\, e^{-\jj 2\pi (f/a) u}\, \dd u
\qquad (u = at) \nonumber\\
&= \frac{1}{a}\, X\!\left(\frac{f}{a}\right).
\label{ch11:scaling-result}
\end{align}
For general real $a \ne 0$ (allowing time reversal),
\[
x(at) \;\longleftrightarrow\; \frac{1}{\abs{a}}\, X\!\left(\frac{f}{a}\right).
\]

If $a > 1$ you compressed the signal in time ($x(at)$ happens faster), and the
spectrum stretched out and shrank in height. If $a < 1$ you stretched the signal
out in time, and the spectrum narrowed and grew taller (Figure~\ref{ch11:fig-scaling}).
Time and frequency pull in opposite directions.

\begin{intuitionbox}
This is the mathematical basis of the uncertainty principle, and you can feel it
without any equations. Clap your hands: a sharp, brief sound (narrow in time) has
a broad, hissy spectrum. Hum a long steady note (wide in time): its spectrum is
a sharp spike. You cannot be brief in time \emph{and} sharp in frequency at the
same time. The scaling theorem is the exact accounting: width-in-time times
width-in-frequency is a constant, so squeezing one widens the other.
\end{intuitionbox}

\begin{goalbox}
The scaling law sets a hard expectation for the wrist controller. Suppose you
want to distinguish a slow wrist roll from a fast finger flick. The flick is
\emph{brief} in time, so by scaling it is \emph{broad} in frequency -- its energy
spreads across many bins. The slow roll is \emph{long} in time and so
\emph{narrow} in frequency. This is good news and a warning at once. Good: the
two gestures live in different frequency regions and are separable. Warning: to
resolve a brief flick you need a short analysis window, which by the very same
law gives you \emph{coarse} frequency resolution. You cannot have fine time
\emph{and} fine frequency resolution simultaneously; Chapter~10 covers
how to choose the trade-off for each gesture.
\end{goalbox}

\begin{figure}[H]
\centering
\begin{tikzpicture}
\begin{axis}[
  width=7cm, height=4.2cm, title={Time domain: $x(t)$ vs $x(2t)$},
  title style={font=\small},
  xlabel={$t$}, xmin=-3, xmax=3, ymin=-0.1, ymax=1.2,
  ytick=\empty, axis lines=left, samples=120, domain=-3:3,
  every axis plot/.append style={thick}, legend style={font=\scriptsize, at={(0.95,0.85)}, anchor=east},
]
\addplot[sigblue] {exp(-x^2)};
\addplot[specred, dashed] {exp(-(2*x)^2)};
\legend{$x(t)$, $x(2t)$}
\end{axis}
\end{tikzpicture}
\hfill
\begin{tikzpicture}
\begin{axis}[
  width=7cm, height=4.2cm, title={Frequency domain: $X(f)$ vs $\frac{1}{2}X(f/2)$},
  title style={font=\small},
  xlabel={$f$}, xmin=-3, xmax=3, ymin=-0.1, ymax=2.0,
  ytick=\empty, axis lines=left, samples=120, domain=-3:3,
  every axis plot/.append style={thick}, legend style={font=\scriptsize, at={(0.95,0.85)}, anchor=east},
]
\addplot[sigblue] {1.77*exp(-(pi*x)^2)};
\addplot[specred, dashed] {0.886*exp(-(pi*x/2)^2)};
\legend{$X(f)$, $\tfrac12 X(f/2)$}
\end{axis}
\end{tikzpicture}
\caption{Scaling. Compressing a Gaussian in time (left, dashed, $x(2t)$) makes
its spectrum wider and shorter (right, dashed). Narrow in time forces broad in
frequency.}
\label{ch11:fig-scaling}
\end{figure}

\section{Duality: the transform is almost its own inverse}
\label{ch11:duality}

Look again at the forward and inverse transforms,
\eqref{ch11:ft-fwd} and \eqref{ch11:ft-inv}. They differ only in the sign of the
exponent. This near-symmetry has a useful consequence: \emph{any} transform
pair, read backwards, is (almost) another transform pair. Precisely, if
$x(t) \longleftrightarrow X(f)$, then
\[
X(t) \;\longleftrightarrow\; x(-f).
\]
That is, take the function that was your \emph{spectrum}, feed it in as a signal,
and its spectrum is your original signal, flipped. (If you use the angular
convention a $2\pi$ appears; in the hertz convention \eqref{ch11:ft-fwd} the
duality is this clean.)

Duality cuts your work in half: every transform pair you derive immediately
gives you a second one for free.

\begin{example}[The rect and the sinc, both directions]
\label{ch11:ex-rect-sinc}
The rectangular pulse of width $T$, $\rect(t/T)$, has transform
\[
\rect(t/T) \;\longleftrightarrow\; T\,\sinc(fT),
\]
where $\sinc(x) = \sin(\pi x)/(\pi x)$. A box in time is a sinc in frequency.
Now apply duality: a sinc in \emph{time} must be a box in \emph{frequency}. And
indeed it is -- a perfect sinc-shaped signal has a perfectly flat,
brick-wall lowpass spectrum. This is exactly the ideal lowpass filter, and the
reason an ideal filter is non-causal (the sinc extends to $t = -\infty$): you get
it for free from duality, no second integral required.
\end{example}

\begin{whybox}
Why does duality hold at all? Because the only difference between analysis and
synthesis -- between measuring how much of each phasor is present and adding the
phasors back up -- is which way the phasors spin ($e^{-\jj}$ vs.\ $e^{+\jj}$).
Flip the spin and analysis becomes synthesis. So computing the spectrum of a
spectrum just runs synthesis with the spin flipped, which reconstructs the
original signal, reversed. The symmetry of the formulas reflects the fact that
projecting onto a basis and reassembling from that basis are the
same operation in opposite directions.
\end{whybox}

\section{A small table, now understood not memorized}
\label{ch11:table}

We can now collect the properties. Read each line as a sentence about phasors,
not as a fact to drill.

\begin{center}
\renewcommand{\arraystretch}{1.5}
\begin{tabular}{@{}lll@{}}
\toprule
\textbf{Property} & \textbf{Time} & \textbf{Frequency}\\
\midrule
Linearity      & $a\,x(t)+b\,y(t)$        & $a\,X(f)+b\,Y(f)$\\
Time shift     & $x(t-t_0)$               & $e^{-\jj 2\pi f t_0}\,X(f)$\\
Freq.\ shift   & $e^{\jj 2\pi f_0 t}\,x(t)$ & $X(f-f_0)$\\
Scaling        & $x(at)$                  & $\frac{1}{\abs{a}}X(f/a)$\\
Duality        & $X(t)$                   & $x(-f)$\\
Conjugation    & $x^{*}(t)$               & $X^{*}(-f)$\\
Convolution    & $(x \conv y)(t)$         & $X(f)\,Y(f)$\\
Multiplication & $x(t)\,y(t)$             & $(X \conv Y)(f)$\\
\bottomrule
\end{tabular}
\end{center}

The conjugation property explains a symmetry that appears in every physical
recording of a real signal. If $x(t)$ is real, then
$x^{*}(t) = x(t)$, so $X^{*}(-f) = X(f)$: the spectrum has
\emph{conjugate symmetry}. Its magnitude is even and its phase is odd. That is
why we usually plot only the positive-frequency half of a real signal's
spectrum -- the negative half is its mirror image and carries no new
information.

The last two lines are the convolution theorem and its dual. The rest of the
chapter is devoted to them.

\begin{historybox}
Joseph Fourier (1768--1830) introduced his series in his 1807 memoir
\emph{On the Propagation of Heat in Solid Bodies}, claiming that \emph{any}
function could be represented as an infinite sum of harmonically related sines and cosines.
The French Academy committee reviewing it -- Lagrange, Laplace, Monge, and Lacroix -- was
skeptical to the point of hostility. Lagrange in particular objected vigorously that smooth,
analytic trigonometric waves could never represent a shape with sharp corners or step
discontinuities. While Fourier's claim that \emph{all} functions could be so represented was an
overreach, his core breakthrough stood. The controversy drove a century of rigorous analysis,
from Dirichlet formulating the first convergence tests in 1829 to Riemann defining modern integration.

For signal processing, these operational properties transformed computation into geometry. Operations
that are cumbersome or intractable in the time domain become simple algebraic rules in frequency:
a time delay becomes a pure linear phase ramp, modulation shifts the spectrum rigidly along the frequency
axis, and temporal convolution collapses into pointwise multiplication.

However, Fourier analysis in its basic form assumes signals extend over all time, from $-\infty$
to $+\infty$. It provides no intrinsic concept of physical causality, memory, or internal system state,
nor does it tell us how to build physical systems that perform these operations step by step.
Chapter~12 (\emph{LTI Systems}) and Chapter~13 (\emph{Convolution}) construct the physical foundations
of causal, time-invariant systems from the ground up, linking Fourier eigenvalues directly to time-domain
impulse responses.
\end{historybox}

\section{Convolution, recalled geometrically}
\label{ch11:convolution-recall}

We will treat convolution in depth in Chapter~13, but we need its picture
now. The convolution of $x$ and $h$ is
\[
(x \conv h)(t) = \int_{-\infty}^{\infty} x(\tau)\, h(t - \tau)\, \dd \tau.
\]
The mechanical recipe: flip $h$, slide it to position $t$, multiply it
point-by-point against $x$, and add up (integrate) the overlap
(Figure~\ref{ch11:fig-conv-slide}). The result at
$t$ is ``how much of $x$ is seen through a flipped-and-shifted copy of $h$.''

\begin{happeningbox}
Forget the flip for a moment and think of what convolution is \emph{for}. $h$ is
a little shape -- an impulse response, a smear, a blur kernel. Convolving $x$
with $h$ stamps a copy of $h$ at every point of $x$, scaled by the value of $x$
there, and adds all the stamps up. A spike in $x$ becomes a copy of $h$; a
smooth ramp in $x$ becomes a smeared ramp. Convolution is \emph{smearing $x$ with
the shape $h$}. Filtering, blurring, echo, averaging, the response of any linear
time-invariant system -- all of it is one signal smeared by another.
\end{happeningbox}

\begin{figure}[H]
\centering
\begin{tikzpicture}
\begin{axis}[
  width=13cm, height=4.5cm,
  xlabel={$\tau$}, xmin=-0.5, xmax=6, ymin=-0.1, ymax=1.3,
  ytick=\empty, axis lines=left, samples=100, domain=-0.5:6,
  every axis plot/.append style={thick},
]
\addplot[sigblue] {exp(-8*(x-1.2)^2) + 0.7*exp(-8*(x-3.2)^2)};
\addplot[specred, dashed] {0.9*exp(-25*(x-4)^2)};
\draw[->, gray] (axis cs:4,1.05) -- (axis cs:4.7,1.05)
  node[right,black]{\small slide $h(t-\tau)$};
\node[sigblue] at (axis cs:1.2,1.15) {\small $x(\tau)$};
\node[specred] at (axis cs:3.4,1.0) {\small $h(t-\tau)$};
\end{axis}
\end{tikzpicture}
\caption{The sliding picture of convolution. The flipped kernel $h(t-\tau)$
(red, dashed) slides along $x(\tau)$ (blue). At each position $t$ the output is
the area of the product -- the overlap. Convolution smears $x$ by the shape of
$h$.}
\label{ch11:fig-conv-slide}
\end{figure}

Computing this overlap integral at every $t$ is expensive. For two length-$N$
discrete signals it costs on the order of $N^2$ multiply-adds: for each of $N$
output positions you sum $N$ products. For a one-second window of IMU data at a
modest rate this is already thousands of operations \emph{per output sample}.
The convolution theorem offers a much cheaper route.

\section{The convolution theorem}
\label{ch11:conv-theorem}

\begin{theorem}[Convolution theorem]
\label{ch11:thm-conv}
If $x(t) \longleftrightarrow X(f)$ and $h(t) \longleftrightarrow H(f)$, then
\[
(x \conv h)(t) \;\longleftrightarrow\; X(f)\,H(f).
\]
Convolution in the time domain is \emph{pointwise multiplication} in the
frequency domain.
\end{theorem}

\begin{proof}
Take the transform of the convolution and swap the order of integration (both
integrals are over all of $\R$, so this is legal when the signals are
well-behaved, e.g.\ absolutely integrable):
\begin{align}
\FT\{x \conv h\}(f)
&= \int_{-\infty}^{\infty}\!\left[\int_{-\infty}^{\infty} x(\tau)\,h(t-\tau)\,\dd\tau\right]
   e^{-\jj 2\pi f t}\, \dd t \nonumber\\
&= \int_{-\infty}^{\infty} x(\tau)
   \left[\int_{-\infty}^{\infty} h(t-\tau)\, e^{-\jj 2\pi f t}\, \dd t\right] \dd\tau.
\label{ch11:proof-swap}
\end{align}
The inner integral is the transform of a \emph{shifted} $h$, which by the
time-shift property \eqref{ch11:time-shift-result} is $e^{-\jj 2\pi f \tau}H(f)$.
Substitute:
\begin{align}
&= \int_{-\infty}^{\infty} x(\tau)\, e^{-\jj 2\pi f \tau}\, H(f)\, \dd\tau
\nonumber\\
&= H(f) \int_{-\infty}^{\infty} x(\tau)\, e^{-\jj 2\pi f \tau}\, \dd\tau
= H(f)\, X(f).
\end{align}
\end{proof}

The proof follows directly from the time-shift property: delaying a kernel by $\tau$ simply applies a phase ramp $e^{-\jj 2\pi f \tau}$ across frequency, which integrates against $x(\tau)$ to produce the Fourier transform of the input.

\begin{whybox}
\textbf{Why convolution is multiplication -- the eigenfunction view.} Feed a
single pure phasor $e^{\jj 2\pi f t}$ into a convolution with $h$:
\[
(e^{\jj 2\pi f \cdot} \conv h)(t)
= \int h(\tau)\, e^{\jj 2\pi f (t-\tau)}\,\dd\tau
= e^{\jj 2\pi f t}\underbrace{\int h(\tau)\, e^{-\jj 2\pi f \tau}\,\dd\tau}_{H(f)}
= H(f)\, e^{\jj 2\pi f t}.
\]
The phasor comes out \emph{unchanged in frequency}, only \emph{scaled} by the
single complex number $H(f)$. Phasors are the \emph{eigenfunctions} of
convolution, and $H(f)$ is the eigenvalue at frequency $f$. Now recall what the
Fourier transform does: it writes $x$ as a sum (integral) of exactly these
phasors. Convolution acts on each phasor independently, multiplying it by its
own eigenvalue $H(f)$. So in the phasor coordinates -- the frequency domain --
convolution is just ``scale each coordinate by $H(f)$,'' which is pointwise
multiplication. In the basis that diagonalizes it, the overlap integral becomes a
simple product.
\end{whybox}

\begin{connectionbox}
This is the same phenomenon as diagonalizing a matrix in linear algebra. A
convolution is a linear, shift-invariant operator; the phasors $e^{\jj 2\pi f t}$
are its eigenvectors; $H(f)$ is the spectrum of eigenvalues. The Fourier
transform is the change of basis that diagonalizes \emph{every} shift-invariant
operator at once. In the DFT, this is literal: a circular convolution is
multiplication by a circulant matrix, and \emph{every} circulant matrix is
diagonalized by the DFT matrix $F$. The discrete convolution theorem is the eigenvalue
decomposition of circulant matrices expressed in frequency-domain notation.
\end{connectionbox}

\begin{whybox}
\textbf{Why this fact is so useful in DSP.} There are three practical reasons:
\begin{enumerate}
\item \emph{It turns design into shaping.} Designing a filter \emph{is} choosing
$H(f)$ -- a curve that is near $1$ where you want to keep signal and near $0$
where you want to kill it. Convolution then realizes that curve. You think in
the picture you actually care about (which frequencies survive) instead of in the
opaque time-domain kernel.
\item \emph{It makes filtering fast.} The overlap integral costs $O(N^2)$.
Multiplication in frequency costs $O(N)$. With the FFT to get in and out of the
frequency domain at $O(N\log N)$, the whole filtering job drops from $O(N^2)$ to
$O(N\log N)$. For a length-$1024$ block that is roughly a hundredfold speedup;
for longer ones, far more.
\item \emph{It unifies the field.} Filtering, correlation, matched detection,
deconvolution (dividing by $H(f)$ to undo a known smear), the response of any LTI
system, even multiplication of polynomials and big integers -- all are
convolutions, and all become multiplications.
\end{enumerate}
\end{whybox}

\subsection{The dual: multiplication in time is convolution in frequency}

By duality (Section~\ref{ch11:duality}), the convolution theorem has a twin:
\[
x(t)\,y(t) \;\longleftrightarrow\; (X \conv Y)(f).
\]
Multiplying two signals in time \emph{convolves} their spectra. We already met a
special case: modulation by a cosine (Section~\ref{ch11:freq-shift}) multiplied
$x(t)$ by $\cos(2\pi f_0 t)$, and the result was the spectrum copied to $\pm
f_0$. That is exactly convolution of $X(f)$ with the two spikes that make up the
cosine's spectrum. The whole modulation picture is one instance of the dual
theorem.

\begin{trapbox}
This dual is also where \emph{windowing} comes from, and it is a very
common source of confusion in spectral work. Whenever you analyze a finite chunk
of a signal you have implicitly multiplied it by a rectangular window in time.
By the dual theorem, your computed spectrum is the true spectrum
\emph{convolved} with the window's transform -- a sinc. That convolution smears
every sharp spectral line into a sinc shape with ringing side-lobes
(``spectral leakage''). You did not measure the true spectrum; you measured it
smeared by the window. Choosing a gentler window (Hann, Hamming) is choosing a
narrower, lower-side-lobe smearing kernel. Spectrogram analysis (Chapter~10) and
spectral estimation (Chapter~21) rely directly on this trade-off.
\end{trapbox}

\begin{figure}[H]
\centering
\begin{tikzpicture}
\begin{axis}[
  width=7cm, height=4.2cm,
  title={Time domain: tone $\times$ window},
  title style={font=\small},
  xlabel={$t$}, ylabel={$x(t) \cdot w(t)$},
  xmin=-2, xmax=2, ymin=-1.3, ymax=1.3,
  axis lines=left, samples=150, domain=-2:2,
  every axis plot/.append style={thick},
]
\addplot[gray!60, dashed] coordinates {(-2,0) (-1,0) (-1,1) (1,1) (1,0) (2,0)};
\addplot[sigblue, domain=-1:1, samples=120] {cos(deg(8*pi*x))};
\node[gray!70!black, font=\scriptsize] at (axis cs:-1.4, 0.6) {window $w(t)$};
\draw[->, gray!70!black, thin] (axis cs:-1.3, 0.45) -- (axis cs:-1.02, 0.2);
\node[sigblue, font=\scriptsize] at (axis cs:1.45, 0.6) {tone};
\draw[->, sigblue, thin] (axis cs:1.3, 0.45) -- (axis cs:1.02, 0.2);
\end{axis}
\end{tikzpicture}
\hfill
\begin{tikzpicture}
\begin{axis}[
  width=7cm, height=4.2cm,
  title={Frequency domain: sinc leakage at $f_0$},
  title style={font=\small},
  xlabel={$f$}, ylabel={$\abs{X \conv W}$},
  xmin=0, xmax=8, ymin=0, ymax=1.2,
  xtick={0,2,4,6,8}, xticklabels={0,2,$f_0$,6,8},
  ytick=\empty, axis lines=left, samples=200, domain=0:8,
  every axis plot/.append style={thick},
]
\addplot[specred] {abs(sin(deg(pi*(x-4)*2)+0.0001)/(pi*(x-4)*2+0.0001))};
\node[specred] at (axis cs:4,1.08) {\footnotesize mainlobe};
\node[quotientorange] at (axis cs:6.2,0.4) {\footnotesize sidelobes};
\draw[->, quotientorange, thick] (axis cs:5.6,0.35) -- (axis cs:4.7,0.22);
\end{axis}
\end{tikzpicture}
\caption{Multiplication in time is convolution in frequency. Truncating a pure tone with a rectangular window of length $T$ (left) convolves its single spectral line at $f_0$ with the window's transform $W(f) = T\,\sinc(fT)$ (right). The sharp impulse broadens into a mainlobe of width $2/T$ flanked by decaying oscillatory sidelobes -- the signature of spectral leakage.}
\label{ch11:fig-window-leakage}
\end{figure}

\section{Fast filtering: the convolution theorem at work}
\label{ch11:fast-filtering}

Let us turn the theorem into an algorithm. To filter $x$ with kernel $h$:
\begin{enumerate}
\item Transform: $X = \FT\{x\}$, $H = \FT\{h\}$ (in practice, FFT).
\item Multiply pointwise: $Y[k] = X[k]\,H[k]$.
\item Inverse transform: $y = \FT^{-1}\{Y\}$.
\end{enumerate}

\begin{lstlisting}[language=Python, caption={Fast filtering via the convolution theorem.}]
import numpy as np

def fast_filter(x, h):
    """Filter signal x with kernel h using the convolution theorem.
    Returns the full linear convolution (length len(x)+len(h)-1)."""
    n = len(x) + len(h) - 1
    nfft = 1 << (n - 1).bit_length()       # next power of two -> fast FFT
    X = np.fft.rfft(x, nfft)               # real FFT, O(n log n)
    H = np.fft.rfft(h, nfft)
    y = np.fft.irfft(X * H, nfft)[:n]      # multiply in freq, then invert
    return y

# sanity check against direct (slow) convolution
x = np.random.randn(2000)
h = np.array([0.2, 0.2, 0.2, 0.2, 0.2])    # 5-point moving average
assert np.allclose(fast_filter(x, h), np.convolve(x, h))
\end{lstlisting}

\begin{trapbox}
\textbf{Circular vs.\ linear convolution -- the trap that gets everyone once.}
The DFT does \emph{not} compute the ordinary (linear) convolution. It computes
\emph{circular} convolution: it treats the signal as one period of a periodic
signal, so the kernel that runs off the right edge wraps around and corrupts the
left edge. The cure is zero-padding. Pad both signals out to length at least
$\text{len}(x) + \text{len}(h) - 1$ before the FFT, and the wrap-around lands
only in the padding, where it does no harm. The code above does exactly this --
the line \texttt{nfft} chooses a length large enough that the circular result
equals the linear one. Forget the pad and you will see mysterious garbage at the
start of your filtered signal.
\end{trapbox}

\begin{intuitionbox}
Is the FFT route always faster? No -- there is a crossover. For a very short
kernel (say a 5-tap average) direct convolution is cheaper per output sample than
two FFTs and an inverse. The frequency route wins decisively once the kernel
gets long (tens to hundreds of taps) or you must filter a long stream. For
streaming, the standard trick is \emph{overlap-add} (or overlap-save): chop the
long input into blocks, FFT-filter each block, and stitch the overlapping tails
back together. That keeps the $O(N\log N)$ advantage while bounding latency and
memory -- exactly what you need on an embedded target.
\end{intuitionbox}

\begin{buildbox}
On a microcontroller at the wrist you rarely FFT-filter the raw IMU stream;
sample rates are low (say 100--500\,Hz) and kernels are short, so a direct FIR
loop wins and has predictable, jitter-free timing. Where the FFT pays off on the wrist
is in \emph{feature extraction}: every analysis frame you take a
real FFT of a short window (CMSIS-DSP's \texttt{arm\_rfft\_fast\_f32} on a
Cortex-M4F does a 256-point real FFT in well under a millisecond) and read band
energies straight off $\abs{X[k]}^2$. Rule of thumb for the build: \emph{filter}
the time stream with a short FIR (cheap, low latency), and use the \emph{FFT} to
\emph{describe} each window for the classifier. The two jobs call for
different tools, and this chapter justifies both.
\end{buildbox}

\section{Correlation: convolution's close cousin}
\label{ch11:correlation}

Correlation answers a different but related question: \emph{how much does $x$ look
like $y$ when $y$ is slid to lag $\tau$?} The cross-correlation is
\[
r_{xy}(\tau) = \int_{-\infty}^{\infty} x^{*}(t)\, y(t + \tau)\, \dd t.
\]
It is convolution \emph{without the flip} (with a conjugate on $x$). The flip is
precisely the difference: convolution flips one signal, correlation does not. So
$r_{xy}(\tau) = (x^{*}(-\cdot) \conv y)(\tau)$, and the convolution theorem
immediately gives the \emph{correlation theorem}:
\[
r_{xy}(\tau) \;\longleftrightarrow\; X^{*}(f)\,Y(f).
\]
Cross-correlate in time $=$ multiply by the conjugate spectrum in frequency. As
with filtering, this makes correlation fast: FFT both signals, multiply one by
the conjugate of the other, inverse FFT.

\begin{intuitionbox}
Why a conjugate and not a plain product? Because correlation measures
\emph{alignment}. In the frequency domain, $X^{*}(f)\,Y(f)$ has magnitude
$\abs{X}\abs{Y}$ (how much energy both share at frequency $f$) and phase
$\angle Y - \angle X$ (how far out of step they are at $f$). When $x$ and $y$ are
the same signal shifted by $\tau_0$, the time-shift property says their phases
differ by exactly the ramp $-2\pi f \tau_0$; multiplying by the conjugate
\emph{cancels} their common shape and leaves precisely that ramp, whose inverse
transform is a spike at $\tau_0$. Correlation finds the lag by letting all
frequencies vote, and the conjugate is what makes the votes line up.
\end{intuitionbox}

The special case $y = x$ gives the \emph{autocorrelation} $r_{xx}(\tau)$, whose
transform is $X^{*}(f)X(f) = \abs{X(f)}^2$ -- the \emph{power spectrum}. This is
the Wiener--Khinchin theorem in miniature: the autocorrelation and the power
spectrum are a Fourier pair. It underlies pitch detection, spectral
estimation, and measuring how ``rhythmic'' a noisy signal is.

\begin{connectionbox}
For a random signal, the autocorrelation $r_{xx}(\tau) = \E[x(t)\,x(t+\tau)]$ is
an \emph{expectation} (probability companion), and its Fourier transform is the
power spectral density. This is the bridge between this chapter's deterministic
algebra and the statistics of noise: the convolution/correlation machinery
carries over verbatim to random signals, with integrals reinterpreted as
expectations. We build this out in Chapter~20.
\end{connectionbox}

\begin{applicationbox}
\textbf{Matched filtering.} To detect a known template $s(t)$ buried in noise,
correlate the incoming signal against $s$. The output peaks at the lag where the
template best aligns -- and among all linear filters this maximizes the
signal-to-noise ratio at the detection instant (provably optimal for white
noise). GPS receivers correlate against known pseudo-random codes; radar
correlates the echo against the transmitted chirp; and a gesture recognizer can
correlate a live IMU window against a stored ``tap'' template to flag candidate
taps. By the correlation theorem, every one of these is a multiply-by-conjugate
in the frequency domain, hence FFT-fast.
\end{applicationbox}

\begin{goalbox}
For the wrist controller, correlation is how you turn ``does this window contain
gesture $g$?'' into a number. Record a clean prototype of each coarse gesture
(double-tap, wrist twist, snap) and, frame by frame, compute the normalized
cross-correlation peak between the live window and each prototype. A high,
well-localized peak is a detection; its lag tells you \emph{when}. This is robust
and cheap, and it degrades gracefully -- exactly the property you want on a
device that must never fire a wrong note in the middle of a performance.

Know its ceiling, though. Template correlation on a wrist-worn IMU
detects \emph{coarse, whole-hand} motions well. It cannot recover which
individual finger moved or by how much: the IMU sits proximal to the fingers and
simply does not see independent finger joint angles. Distinguishing finger-level
events needs electromyography (EMG) or force myography (FMG) sensing the
forearm muscles, and \emph{continuous} finger pose needs sensors on the fingers
themselves (rings, gloves). Correlation makes the gestures the IMU \emph{can} see
reliable, but it cannot recover information the sensor never captured.
\end{goalbox}

\section{Worked examples}
\label{ch11:worked}

\begin{example}[Delaying a pulse]
\label{ch11:ex-delay}
A rectangular pulse $p(t) = \rect(t/T)$ has spectrum $P(f) = T\,\sinc(fT)$,
which is real (the pulse is centered at $0$, hence even). Now delay it by
$t_0 = T/2$ so it starts at $0$ and ends at $T$. By the time-shift property the
new spectrum is
\[
e^{-\jj 2\pi f (T/2)}\, T\,\sinc(fT).
\]
The \emph{magnitude} $\abs{T\,\sinc(fT)}$ did not move at all -- delaying a
signal never changes which frequencies are present or how strong they are. What
changed is the phase, which now carries a $-\pi f T$ ramp encoding ``this pulse
arrives $T/2$ seconds late.'' All timing information in a signal lives in the
\emph{phase} of its spectrum; the magnitude is timing-blind.
\end{example}

\begin{example}[Convolving two boxes gives a triangle]
\label{ch11:ex-box-box}
Convolve a width-$T$ box with itself. In time, sliding one box across the other
makes the overlap grow linearly, peak at full overlap, then shrink -- a
\emph{triangle} of width $2T$. Now check it through the convolution theorem.
Each box transforms to $T\,\sinc(fT)$, so the convolution transforms to
$[T\,\sinc(fT)]^2 = T^2\,\sinc^2(fT)$. And indeed the known transform of a
triangle is $\sinc^2$. Two ways to the same place: slide-and-overlap in time, or
square the sinc in frequency. The frequency route also explains a feature you
might not have predicted -- because $\sinc^2 \ge 0$ everywhere, the triangle's
spectrum has \emph{no sign changes}, unlike the box's $\sinc$, which dips
negative. Squaring in frequency lifted those negative lobes up. This is exactly
why a triangular (Bartlett) window leaks less than a rectangular one: its
spectral side-lobes are the \emph{square} of the rectangle's, hence much smaller.
\end{example}

\begin{example}[Smoothing as lowpass filtering]
\label{ch11:ex-smoothing}
A 5-point moving average has kernel
$h[n] = \tfrac15(\delta[n] + \delta[n{-}1] + \dots + \delta[n{-}4])$. Its DTFT is
a sum of phasors -- a Dirichlet kernel -- whose magnitude is
\[
\abs{H(f)} = \frac{1}{5}\left|\frac{\sin(5\pi f)}{\sin(\pi f)}\right|,
\]
with $f$ in cycles per sample. It equals $1$ at $f = 0$ (DC passes untouched,
since the weights sum to $1$) and falls off toward higher $f$, with nulls where
$5\pi f$ is a nonzero multiple of $\pi$ (Figure~\ref{ch11:fig-movavg}). So convolving with this kernel
\emph{multiplies} the input spectrum by a lowpass curve: slow trends survive,
fast wiggles are attenuated, and any frequency landing exactly on a null is
annihilated. ``Averaging smooths the data'' and ``averaging is a lowpass filter''
are the same sentence read in the two domains -- and the convolution theorem is
the dictionary between them. This is precisely the smoother you would run on a
jittery wrist-IMU axis before extracting features.
\end{example}

\begin{figure}[H]
\centering
\begin{tikzpicture}
\begin{axis}[
  width=12cm, height=4.6cm,
  xlabel={frequency $f$ (cycles/sample)}, ylabel={$\abs{H(f)}$},
  xmin=0, xmax=0.5, ymin=0, ymax=1.1,
  xtick={0,0.1,0.2,0.3,0.4,0.5}, ytick={0,0.5,1},
  axis lines=left, samples=200, domain=0.0005:0.5,
  every axis plot/.append style={thick}, clip=true,
]
\addplot[sigblue] {abs(sin(deg(5*pi*x))/sin(deg(pi*x)))/5};
\node[sigblue] at (axis cs:0.07,1.0) {\small DC passes};
\node[specred] at (axis cs:0.25,0.18) {\small null at $f=1/5$};
\addplot[specred, only marks, mark=*, mark size=1.4pt] coordinates {(0.2,0) (0.4,0)};
\end{axis}
\end{tikzpicture}
\caption{Magnitude response of the 5-point moving average,
$\abs{H(f)} = \tfrac15\abs{\sin(5\pi f)/\sin(\pi f)}$. It equals $1$ at DC,
falls off toward higher frequencies, and has nulls at $f = 1/5$ and $f = 2/5$
cycles/sample. Averaging is a lowpass filter.}
\label{ch11:fig-movavg}
\end{figure}

\begin{walkbox}
Suppose you compute the cross-correlation of two recordings of the \emph{same}
gesture taken a fraction of a second apart, and the correlation peak sits at lag
$\tau = +12$ samples. What does that tell you, and what would you read off the
\emph{phase} of $X^{*}(f)Y(f)$ to confirm it? (Answer: the second recording lags
the first by 12 samples; the phase should be the linear ramp $-2\pi f \cdot 12$
across frequency -- timing lives in phase slope, and a clean straight ramp is the
signature of a pure delay between two otherwise identical signals.)
\end{walkbox}

\begin{connectionbox}
\textbf{Capabilities, Downsides, and Forward Roadmap}
\begin{itemize}
  \item \textbf{What this made possible}: The operational properties of the Fourier transform convert difficult
  time-domain calculus into elementary frequency-domain algebra. Time delays become linear phase ramps,
  scaling trades duration against bandwidth, and temporal convolution collapses into pointwise multiplication.
  \item \textbf{Limitations}: The convolution theorem is an algebraic relationship, but it
  says nothing about the physical structure of devices that process signals in time. Real hardware cannot wait
  for future samples ($t \to +\infty$), operates sequentially sample-by-sample, and requires strict guarantees
  on causality, memory, and stability.
  \item \textbf{Where we go next}: Chapter 12 builds the physical framework of \textbf{Linear Time-Invariant (LTI) Systems}.
  We discover why linearity and time-invariance are the exact physical conditions that make complex sinusoids
  eigenfunctions, and Chapter 13 proves from first principles why every such system is completely
  governed by the convolution sum.
\end{itemize}
\end{connectionbox}

\section{Exercises}
\label{ch11:exercises}

\begin{exercise}
Using only the time-shift and linearity properties, find the Fourier transform of
$x(t) = \delta(t - 2) + \delta(t + 2)$. Show that the result is real and express
it as a single cosine in $f$. Explain geometrically why a symmetric pair of
impulses gives a purely real (cosine) spectrum.
\end{exercise}

\begin{exercise}
A signal $x(t)$ has transform $X(f)$. Find the transform of
$y(t) = x(t)\cos(2\pi f_0 t)\allowbreak\cos(2\pi f_0 t)$ in terms of $X$. (Hint: apply the
modulation property twice, or use a double-angle identity first.) How many copies
of $X$ appear, and at what frequencies and heights?
\end{exercise}

\begin{exercise}
Prove the scaling property for $a < 0$ directly from the definition, being
careful with the limits of integration when you substitute $u = at$, and confirm
that the $1/\abs{a}$ (not $1/a$) factor is correct. What does $a = -1$ do to a
signal and to its spectrum?
\end{exercise}

\begin{exercise}
Starting from the pair $\rect(t/T) \leftrightarrow T\,\sinc(fT)$, use duality to
write the transform of $\sinc(t/B)$. Sketch its magnitude. Why does this make
$\sinc$ the ideal lowpass filter, and why can such a filter never run in real
time?
\end{exercise}

\begin{exercise}
Let $h$ be a 3-point kernel $[1, -2, 1]$. Compute $H(f)$ (DTFT) and identify the
filter type from its magnitude. Where are its zeros, and what does the kernel do
to a constant (DC) input and to a slowly varying ramp? Relate your answer to the
second derivative.
\end{exercise}

\begin{exercise}
You filter a length-$1000$ signal with a length-$50$ kernel using the FFT, but
you set the FFT length to exactly $1000$ instead of zero-padding. Describe
precisely what corruption appears in the output and where. How long must the FFT
be to avoid it, and why?
\end{exercise}

\begin{exercise}
Show that the autocorrelation of any real signal is even,
$r_{xx}(-\tau) = r_{xx}(\tau)$, and is maximized at $\tau = 0$. Then show, using
the correlation theorem, that its Fourier transform $\abs{X(f)}^2$ is real and
nonnegative. Why must a valid power spectrum be nonnegative?
\end{exercise}

\begin{exercise}
Estimate the operation counts for filtering a length-$N$ signal with a length-$M$
kernel (a) directly and (b) via FFT (one $N$-point FFT each for signal and
kernel, $N$ complex multiplies, one inverse FFT). For $N = 4096$, find the kernel
length $M$ at which the FFT method becomes cheaper. State your assumptions.
\end{exercise}

\begin{exercise}
A wrist IMU axis is contaminated by a 50\,Hz tone from nearby mains wiring,
riding on top of the gesture signal. In the frequency domain, what is the
cleanest way to remove it, and what time-domain operation does that correspond
to? What artifact does an abrupt ``notch'' (setting one bin to zero) create in
the time domain, and how does the convolution theorem predict it?
\end{exercise}

\begin{exercise}
Two microphones (or two IMUs) record the same event with an unknown relative
delay. Explain step by step how to recover the delay using one FFT per channel,
a conjugate multiply, and one inverse FFT, and how to read the delay off the
result. What goes wrong if the two recordings are not actually shifted copies of
each other but merely similar?
\end{exercise}

\begin{exercise}
(Open-ended.) Argue both sides of the claim ``the convolution theorem is the most
useful fact in DSP.'' List three concrete tasks it makes possible or fast, and
one situation where reaching for it is the wrong move (where direct time-domain
processing is simpler or faster). Tie your wrong-move example to the wrist
controller.
\end{exercise}
